\subsection{RATIONAL LINEAR FORMS}

Let V be a right vector space over K with a \(K^{\prime}\) -structure \(V^{\prime}\) . As \(K_{d}^{\prime}\) is a \(K^{\prime}\) -structure on the right vector K-space \(K_{d}\) , we may define linear forms \(x^{*} \in V^{*}\) , rational over \(K^{\prime}\) , as the linear mappings of V into \(K_{d}\) , rational over \(K^{\prime}\) for the

K'-structures on V and \(K_{d}\) . By virtue of no. 3, Proposition 3, the set \(R'\) of these linear forms is the image of the dual \(V'^*\) of \(V'\) under the composite mapping

\[
\mathrm{V} ^ {\prime *} \xrightarrow {\phi} \mathrm{K} \otimes_ {\mathrm{K} ^ {\prime}} \mathrm{V} ^ {\prime *} \xrightarrow {\upsilon} \mathrm{V} ^ {*}\tag{1}
\]

where \(\phi(x^{\prime*}) = 1 \otimes x^{\prime*}\) and \(\upsilon(\xi \otimes x^{\prime*})\) is the linear form \(y^{*}\) on V such that \(y^{*}(x^{\prime}) = \xi\langle x^{\prime*}, x^{\prime} \rangle\) for all \(x^{\prime} \in V^{\prime}\) (§ 5, no. 4). We know that this mapping is injective (§ 7, no. 9, Propositions 20 and 19) and clearly \(R^{\prime}\) is a left vector sub- \(K^{\prime}\) -space of \(V^{*}\) ; moreover every subset of \(R^{\prime}\) free over \(K^{\prime}\) is free over K. But in general \(R^{\prime}\) does not necessarily generate \(V^{*}\) over K and does not therefore define a \(K^{\prime}\) -structure on \(V^{*}\) (Exercise 2). However, if V is of finite dimension \(n\) over K, \(V^{\prime*}\) is of dimension \(n\) over \(K^{\prime}\) and \(R^{\prime}\) then defines canonically a \(K^{\prime}\) -structure on \(V^{*}\) .

PROPOSITION 5. Let V be a right vector space over K, V' a K'-structure on V and W a vector sub-K-space of V. For W to be rational over K', it is necessary and sufficient that there exist a set H ⊂ V* of rational linear forms over K' such that W is the orthogonal of H in V (§ 2, no. 4).

Let H be a subset of V* whose elements are rational linear forms over K'. For all \(x^{*} \in H\) , the kernel of \(x^{*}\) is a vector sub-K-space of V, rational over K' (no. 3, Corollary 2 to Proposition 3); the intersection of these kernels is therefore also a vector sub-K-space of V, rational over K' (no. 2, Corollary 1 to Proposition 2).

Conversely, let W be a vector sub-K-space of V rational over \(K'\) , so that W is identified with \(W' \otimes_{K'} K\) , where \(W' = W \cap V'\) (no. 2, Proposition 2). For a linear form \(x'^* \in V'^*\) to be zero on \(W'\) , it is necessary and sufficient that the linear form \(x^* \in V^*\) which corresponds to it under (1) be zero on W, for by no. 3, Corollary 1 to Proposition 3,

\[
\operatorname{Ker} (x ^ {*}) = \left(\operatorname{Ker} (x ^ {\prime *})\right) \otimes_ {\mathbb {K} ^ {\prime}} \mathrm{K} \quad \text { and } \quad \operatorname{Ker} (x ^ {* ^ {\prime}}) = (\operatorname{Ker} (x ^ {*})) \cap \mathrm{V} ^ {\prime}.
\]

Let \(H'\) be the orthogonal of \(W'\) in \(V'^*\); we know (§ 7, no. 5, Theorem 7) that \(W'\) is the orthogonal of \(H'\) in \(V'\) ; if \(H\) is the image of \(H'\) in \(V*\) under the mapping (1), it follows from the above that \(W\) is the orthogonal of \(H\) in \(V\) , taking account of no. 7, Corollary to Proposition 14.

